\chapter{数え上げ}
    \section{\texorpdfstring{`for $i=1, \dotsc, n(\in \naturalNumbers)$'$\iff$`for $d \mid n, \text{ for } i=1,\dotsc,n \text{ where } \gcd(i,n)=d$'}{数え上げと最大公約数}}
        \begin{shadebox}
            次は同値である。
            \begin{enumerate}
                \item `for $i=1, \dotsc, n(\in \naturalNumbers)$'
                \item `for $d \mid n, \text{ for } i=1,\dotsc,n \text{ where } \gcd(i,n)=d$'
                \item `for $d \mid n, \text{ for } j=1,\dotsc,n/d \text{ where } \gcd(j,n/d)=1$'
            \end{enumerate}
        \end{shadebox}
        1から$n$までの自然数に対して，それと$n$との最大公約数は$n$の約数のどれかと一致するから，最初の$\iff$が成り立つ。
        2つ目の$\iff$については$i=dj$なる変数変換を施すことで分かる。
        次のような使い方が考えられる。
        \begin{itemize}
            \item $\sum_{i=1}^n f(i) = \sum_{d \mid n} \sum_{\genfrac{}{}{0pt}{}{i=1, \dotsc, n}{\gcd(i,n)=d}} f(i) = \sum_{d \mid n} \sum_{\genfrac{}{}{0pt}{}{j=1, \dotsc, n/d}{\gcd(j,n/d)=1}} f(dj)$
            \item $\prod_{i=1}^n f(i) = \prod_{d \mid n} \prod_{\genfrac{}{}{0pt}{}{i=1, \dotsc, n}{\gcd(i,n)=d}} f(i) = \prod_{d \mid n} \prod_{\genfrac{}{}{0pt}{}{j=1, \dotsc, n/d}{\gcd(j,n/d)=1}} f(dj)$
        \end{itemize}
    \section{包除原理}
        \subsection{補題: \texorpdfstring{$A\subseteq U \Rightarrow |A| = \sum_{x\in U}{|A\cap\{x\}|}$}{A⊆Uのときの|A|の計算}}
            \label{補題: 要素数のカウント。親集合の各要素について子集合に含まれるかどうか調べる。}
            \begin{commentary}
                \quad\par
                当たり前すぎて証明に困る。$U$の各要素について，それが$A$に含まれるかどうかを虱潰しに調べているだけ。
            \end{commentary}
            \begin{shadebox}
                任意の集合$A_1,\dotsc,A_n$に対して次の等式が成り立つ。
                \begin{align*}
                    |A_1\cup\dotsb\cup A_n| &= |A_1|+\dotsb +|A_n| \\
                    &- |A_1\cap A_2| - |A_1\cap A_3| - \dotsb - |A_{n-1}\cap A_n| \\
                    &+ |A_1\cap A_2\cap A_3| + |A_1\cap A_2\cap A_4| + \dotsb + |A_{n-2}\cap A_{n-1}\cap A_n| \\
                    &- \dotsb +(-1)^{n-1}|A_1\cap\dotsb\cap A_n|
                \end{align*}
                すなわち
                \begin{equation}
                    \abs{\bigcup_{i=1}^n A_i} = \sum_{k=1}^n{(-1)^{k-1}\sum_{I\in\combi{\numset}{k}}}{\abs{\bigcap_{i\in I}A_i}}
                \end{equation}
                もっと短く書けば
                \[\abs{\bigcup_{i=1}^n A_i} = \sum_{I\in 2^\numset\setminus\left\{\emptyset\right\}}{(-1)^{|I|-1}\abs{\bigcap_{i\in I}A_i}}\]
            \end{shadebox}
            \begin{proof}
                \quad\par
                式(1)を証明する。
                $U\coloneq\bigcup_{i=1}^n A_i$とする。
                補題\ref{補題: 要素数のカウント。親集合の各要素について子集合に含まれるかどうか調べる。}より
                \begin{align*}
                    \textrm{(1)の左辺} &= \sum_{x\in U}{\abs{\left(\bigcup_{i=1}^n A_i\right)\cap\{x\}}} = \sum_{x\in U}1 \\
                    \textrm{(1)の右辺} &= \sum_{k=1}^n{(-1)^{k-1}\sum_{I\in\combi{\numset}{k}}}{\sum_{x\in U}{\abs{\left(\bigcap_{i\in I}A_i\right)\cap\{x\}}}} \\
                    &= \sum_{k=1}^n{(-1)^{k-1}\sum_{x\in U}{\sum_{I\in\combi{\numset}{k}}{\abs{\left(\bigcap_{i\in I}A_i\right)\cap\{x\}}}}} \quad(\textrm{和の順序交換}) \\
                    &= \sum_{x\in U}{\sum_{k=1}^n{(-1)^{k-1}\sum_{I\in\combi{\numset}{k}}{\abs{\left(\bigcap_{i\in I}A_i\right)\cap\{x\}}}}} \quad(\textrm{和の順序交換}) \\
                    &= \sum_{x\in U}{c(x)} \quad\textrm{where}\quad c(x) \coloneq \sum_{k=1}^n{(-1)^{k-1}\sum_{I\in\combi{\numset}{k}}{\abs{\left(\bigcap_{i\in I}A_i\right)\cap\{x\}}}}
                \end{align*}
                任意の$x\in U$に対して$c(x)=1$であることを示す。
                各$x$に対して，$x$に依存して決まる次のような自然数$m(x)$が存在する。
                \begin{center}
                    「$x$は$m(x)$個の集合$A_{i_1},\dotsc,A_{i_{m(x)}}$に共通して属し，他の$n-m(x)$個の集合には属さない。」
                \end{center}
                従って，$c(x)$の$\sum_{I\in\combi{\numset}{k}}{\abs{\left(\bigcap_{i\in I}A_i\right)\cap\{x\}}}$は$k>m(x)+1$に対して$0$であるから$c(x)$は次のように書き直される。
                \[c(x) = \sum_{k=1}^{\textcolor{red}{m(x)}}{(-1)^{k-1}\sum_{I\in\combi{\numset}{k}}{\abs{\left(\bigcap_{i\in I}A_i\right)\cap\{x\}}}}\]
                上式の$\abs{\left(\bigcap_{i\in I}A_i\right)\cap\{x\}}$が$1$になるのは上述の$m(x)$個の集合から$k$個の集合を選んできた時だけであるから，$c(x)$はさらに次のように書き直される。
                \begin{align*}
                    c(x) &= \sum_{k=1}^{m(x)}{(-1)^{k-1}\sum_{I\in\combi{\{\textcolor{red}{i_1,\dotsc,i_{m(x)}}\}}{k}}{1}} = \sum_{k=1}^{m(x)}{(-1)^{k-1}\combi{m(x)}{k}} \\
                    &= -\sum_{k=1}^{m(x)}{(-1)^k\combi{m(x)}{k}} = -\left(\sum_{k=\textcolor{red}{0}}^{m(x)}{(-1)^k\combi{m(x)}{k}} - 1\right) \\
                    &= -\left((1-1)^{m(x)} - 1\right) = 1
                \end{align*}
            \end{proof}
            \subsection{あの等式をどうやって思いついたか}
                包除原理の等式を証明することはできたが，一番最初にあの等式を思い付いた人は何を考えたのだろうか。
                ここでは，あの等式を導き出す自然なアルゴリズムを考えてみる。
                \par
                まず我々は「メモリ」を\textcolor{red}{1つ}用意して0で初期化する。
                次に，$U\coloneq\bigcup_{i=1}^n A_i$とし，$U$の\textcolor{red}{各}要素に「カウンタ」を取り付けて0で初期化する。
                \par
                まず次の操作を実行する。
                「任意の$A_i$について，$A_i$に属する各要素のカウンタを+1する。カウンタを操作する度にメモリも+1する。この操作を$i=1,\dotsc,n$すべてに対して行う。」
                この操作の実行後，メモリの値は$|A_1|+\dotsb +|A_n|$になっている。
                そしてカウンタの値は次のようになっている。
                \begin{itemize}
                    \item 1つの集合に属しており，2つ以上の集合に属していない要素のカウンタは1
                    \item 2つの集合に属しており，3つ以上の集合に属していない要素のカウンタは2
                    \item $\cdots$
                    \item $n$個の集合に属している要素のカウンタは$n$
                \end{itemize}
                当然だ。複数の集合が被さっている部分の要素はオーバーカウントされている。
                \par
                そこで，「2つの集合に属しており，3つ以上の集合に属していない要素」のオーバーカウントを解消するために次の操作を行う。
                「任意の2つの集合$A_i,A_j$について，$A_i\cap A_j$に属する要素のカウンタを-1する。カウンタを操作する度にメモリも-1する。この操作を全ての$\{i,j\}\in\combi{\numset}{2}$に対して行う。」
                この操作の実行後，メモリの値は$|A_1|+\dotsb +|A_n| - |A_1\cap A_2| - |A_1\cap A_3| - \dotsb - |A_{n-1}\cap A_n|$になっている。
                そしてカウンタの値は次のようになっている。
                \begin{itemize}
                    \item 1つの集合に属しており，2つ以上の集合に属していない要素のカウンタは1
                    \item 2つの集合に属しており，3つ以上の集合に属していない要素のカウンタは$2-\combi{2}{2}=1$
                    \item 3つの集合に属しており，4つ以上の集合に属していない要素のカウンタは$3-\combi{3}{2}=\textcolor{red}{0}$
                    \item 4つの集合に属しており，4つ以上の集合に属していない要素のカウンタは$4-\combi{4}{2}=-2$
                    \item 5つの集合に属しており，4つ以上の集合に属していない要素のカウンタは$5-\combi{5}{2}=-5$
                    \item $\cdots$
                    \item $n$個の集合に属している要素のカウンタは$n - \combi{n}{2} < 0$
                \end{itemize}
                (※ $\combi{?}{2}$が現れる理由は次のとおり。
                例えば「3つの集合に属しており，4つ以上の集合に属していない要素」は，上述の操作により，自分が属している3つの集合から2つを選んで作られた$\combi{3}{2}$通りの集合のパターンの数だけカウンタが-1される。)
                \par
                やりすぎたようだ。3つ以上の集合に属している要素をアンダーカウントしている。
                \par
                そこで，「3つの集合に属しており，4つ以上の集合に属していない要素」を数え直すために次の操作を実行する。
                「任意の3つの集合$A_i,A_j,A_k$について，$A_i\cap A_j\cap A_k$に属する要素のカウンタを+1する。カウンタを操作する度にメモリも+1する。この操作を全ての$\{i,j,k\}\in\combi{\numset}{3}$に対して行う。」
                この操作の実行後，メモリの値は$|A_1|+\dotsb +|A_n| - |A_1\cap A_2| - |A_1\cap A_3| - \dotsb - |A_{n-1}\cap A_n| + |A_1\cap A_2\cap A_3| + |A_1\cap A_2\cap A_4| + \dotsb + |A_{n-2}\cap A_{n-1}\cap A_n|$になっている。
                そしてカウンタの値は次のようになっている。
                \begin{itemize}
                    \item 1つの集合に属しており，2つ以上の集合に属していない要素のカウンタは1
                    \item 2つの集合に属しており，3つ以上の集合に属していない要素のカウンタは$2-\combi{2}{2}=1$
                    \item 3つの集合に属しており，4つ以上の集合に属していない要素のカウンタは$3-\combi{3}{2}+\combi{3}{3}=1$
                    \item 4つの集合に属しており，4つ以上の集合に属していない要素のカウンタは$4-\combi{4}{2}+\combi{4}{3}=2$
                    \item 5つの集合に属しており，4つ以上の集合に属していない要素のカウンタは$5-\combi{5}{2}+\combi{5}{3}=5$
                    \item $\cdots$
                    \item $n$個の集合に属している要素のカウンタは$n - \combi{n}{2} +\combi{n}{3} > 0$
                \end{itemize}
                4つ以上の集合に含まれる要素はオーバーカウントされている。
                \par
                もう気づいたと思うが，これまでやってきたように「足しすぎては引き，引きすぎては足し，$\cdots$」を繰り返すと最終的にメモリの値は包除原理の主張である$|A_1|+\dotsb +|A_n| - |A_1\cap A_2| - |A_1\cap A_3| - \dotsb - |A_{n-1}\cap A_n| + |A_1\cap A_2\cap A_3| + |A_1\cap A_2\cap A_4| + \dotsb + |A_{n-2}\cap A_{n-1}\cap A_n| - \dotsb +(-1)^{n-1}|A_1\cap\dotsb\cap A_n|$になり，カウンタの値は次のようになる。
                \begin{itemize}
                    \item 1つの集合に属しており，2つ以上の集合に属していない要素のカウンタは1
                    \item 2つの集合に属しており，3つ以上の集合に属していない要素のカウンタは$2-\combi{2}{2}=1$
                    \item 3つの集合に属しており，4つ以上の集合に属していない要素のカウンタは$3-\combi{3}{2}+\combi{3}{3}=1$
                    \item 4つの集合に属しており，4つ以上の集合に属していない要素のカウンタは$4-\combi{4}{2}+\combi{4}{3}-\combi{4}{4}=1$
                    \item 5つの集合に属しており，4つ以上の集合に属していない要素のカウンタは$5-\combi{5}{2}+\combi{5}{3}-\combi{5}{4}+\combi{5}{5}=1$
                    \item $\cdots$
                    \item $n$個の集合に属している要素のカウンタは$n + \sum_{i=2}^n{(-1)^{i-1}\combi{n}{i}} = -\sum_{i=1}^n{(-1)^i\combi{n}{i}} = -\left(\sum_{i=0}^n{(-1)^i\combi{n}{i}} - 1\right) = -\left((1-1)^n - 1\right) = 1$
                \end{itemize}
                全ての要素を過不足無くカウントできていることがわかる。
    \section{配列の飛び飛びマーキング}
        \begin{shadebox}
            $m<n$を互いに素な自然数とする。
            自然数$k$を1から順番に増やしながら，長さ$n$の配列Aの$km\%n$番目の要素をマークしていけば，$k=n$の時にすべての要素を丁度1回ずつマーク完了する。
        \end{shadebox}
        \begin{proof}
            \quad\par
            $m,n$は互いに素であるから，A[$n-1$]が初めてマークされるのは$km=$lcm($m,n$)，すなわち$k=n$の時である。
            仮に，それまでのある時刻でA[0]$\sim$A[$n-2$]のどれかが2回マークされたとすると，動作の周期性から，$k$をどれだけ増やしていってっも，それまでにマークされた要素を何度もマークするだけで，A[$n-1$]はいつまでもマークされない。これは不合理である。
            また，仮に$k=n$の時点で未マークの要素があるとすると，鳩の巣原理から少なくとも1つの要素が2回マークされていることになり，不合理である。
        \end{proof}